\documentclass[11pt]{article}

\usepackage[margin=1in]{geometry}
\usepackage{amsmath, amssymb}
\usepackage{booktabs}
\usepackage[hidelinks]{hyperref}
\usepackage[htt]{hyphenat}
\sloppy

\newcommand{\dt}{\tau}
\newcommand{\uh}{\mathbf u}
\DeclareMathOperator{\cond}{\kappa}

\title{PINO for Allen--Cahn: what the natural generalisation is,\\ and where it stops resembling Poisson}
\date{\today}

\begin{document}
\maketitle

\noindent
Working notes for the Allen--Cahn benchmark, not paper material. The question: given
\texttt{PINOPoissonLoss}, what is the natural PINO for the autoregressive Allen--Cahn setting, and
which of its design choices carry over unchanged? The short answer is that the construction
generalises cleanly --- and is already implemented as
\texttt{tensorpils.baselines.pino.PINOACLoss} --- but \emph{three} of the Poisson version's
choices do not survive, and one of them changes the optimisation geometry substantially.

\section{What PINO does for Poisson}

For $-\Delta u = f$ the reference (\texttt{neuraloperator/physics\_informed}, the Darcy example,
which is the Dirichlet one) forms the residual with second-order central differences on the grid,
evaluates it on interior points only, and reduces with the \emph{relative} $L^p$ ratio
\begin{equation}
  L_{\mathrm{PINO}}(u) \;=\; \frac{\bigl\| -\Delta_h u - f \bigr\|_2}{\|f\|_2},
  \label{eq:poisson}
\end{equation}
averaged over the batch. The boundary condition is imposed \emph{hard, in the model}: the
reference multiplies by $\sin(\pi x)\sin(\pi y)$; we instead zero the boundary nodes, which is the
same operation \texttt{losses.py} applies to our own arms.

\section{The natural generalisation}

Allen--Cahn is not a boundary-value problem but a time-stepper, so the object to make a residual of
is the \emph{step}, not the equation. Our FEM arm builds its residual from exactly the same time
discretisation the reference solver zeros: with $\rho$ the reaction term,
\begin{equation}
  R_{\mathrm{FEM}}(u^{n},u^{n+1}) \;=\; M\frac{u^{n+1}-u^{n}}{\dt} \;+\; a^2 A\,u^{n+1} \;-\; M\rho ,
  \qquad
  \rho = \epsilon^2\bigl(u^{n} - (u^{n+1})^3\bigr)
  \label{eq:fem}
\end{equation}
for the convex--concave (Eyre) split, which is what all our runs use.

The natural PINO is then obtained by the same two substitutions that produce \eqref{eq:poisson}
from the Poisson weak form: \textbf{drop the mass matrix} and \textbf{replace $A$ by
$-\Delta_h$}, the five-point finite-difference Laplacian. This gives
\begin{equation}
  \boxed{\;
  R_{\mathrm{PINO}}(u^{n},u^{n+1}) \;=\; \frac{u^{n+1}-u^{n}}{\dt} \;-\; a^2 \Delta_h u^{n+1}
  \;-\; \epsilon^2\bigl(u^{n} - (u^{n+1})^3\bigr) \; }
  \label{eq:pinoac}
\end{equation}
evaluated pointwise on interior nodes. Nothing here is a choice: given the Poisson construction,
\eqref{eq:pinoac} is forced. It is what \texttt{PINOACLoss} implements.

Over an autoregressive rollout of $R$ steps the loss discount-weights the per-step residuals,
\begin{equation}
  L \;=\; \sum_{k=0}^{R-1} \gamma^{k}\,\bigl\| R_{\mathrm{PINO}}(u^{k},u^{k+1}) \bigr\|^2_{\text{interior}} ,
\end{equation}
matching \texttt{ACGalerkinLoss}'s discount so the two differ only in the residual.

\section{Three things that do not carry over}

\subsection{There is no right-hand side, so the relative reduction is unavailable}

This is the important one. PINO's signature reduction \eqref{eq:poisson} divides by $\|f\|$. The
Allen--Cahn step residual has \emph{no} right-hand side to normalise against --- $u^n$ is the
previous iterate, not a datum --- so the ratio has no analogue and \texttt{PINOACLoss} reduces with
a plain mean square instead.

That is not a cosmetic difference. For Poisson the \texttt{rel} reduction is a \emph{norm}, not a
squared norm: its level sets form a cone, its gradient has constant magnitude $1/\|f\|$ everywhere,
it cannot converge at a fixed step size, and --- the part most often misread --- its normalisation
does \emph{not} improve conditioning (on the tangent space the curvature spread is still
$\cond(A_{\mathrm{FD}})^2$). See \texttt{notes/pino\_loss/pino\_loss\_analysis.tex}. Reverting to
\texttt{mse} puts the Allen--Cahn PINO back in ordinary least-squares geometry, with a Hessian
$\propto J^\top J$ and hence a squared condition number.

\textbf{Consequence for the paper:} the PINO arm is not the same kind of object across the two
PDEs. On Poisson it is a normalised norm; on Allen--Cahn a squared residual. That asymmetry is
forced by the problem, not chosen by us, but it must be stated --- otherwise ``PINO'' looks like
one baseline when it is two.

\subsection{The operator to precondition is not the stiffness matrix}

For Poisson the relevant operator is $A$ itself, with $\cond(A)=\Theta(h^{-2})$. Here the residual
\eqref{eq:fem} is linearised by the step Jacobian
\begin{equation}
  J \;=\; \frac{M}{\dt} \;+\; a^2 A \;+\; 3\epsilon^2 M \operatorname{diag}\bigl((u^{n+1})^2\bigr),
\end{equation}
a \emph{shifted} stiffness matrix. In the convex--concave split the cubic sits at $u^{n+1}$, so the
shift is $c = 1/\dt + 3\epsilon^2$; at the preliminary settings ($\dt=0.01$, $\epsilon=32$) that is
$c = 3172$, which is what \texttt{realistic\_ac} records. The mass shift \emph{helps}: it lifts the
small eigenvalues, so $\cond(J) \ll \cond(A)$ at fixed $h$, and the multigrid preconditioner in
\texttt{tensorpils/preconditioners/multigrid.py} takes exactly this $c$ as a parameter.

None of this touches PINO, which is unpreconditioned by construction --- but it does mean the
Allen--Cahn comparison is \emph{less} favourable to preconditioning a priori than Poisson was,
because the operator being inverted is better conditioned to begin with. Worth knowing before
reading the result.

\subsection{Dropping $M$ means something different here}

For Poisson, discarding the mass matrix amounts to \emph{mass lumping}: the load $Mf$ is replaced
by pointwise $f$, and since $M$'s interior row sums are exactly $h^2$, one finds
$L_{\mathrm{mse}} = (h^{-4}/N)\,\|A_{P_1}u - M_L f\|^2$ --- the finite-difference PINO is our
least-squares loss with $P_1$ triangles instead of $Q_1$ quads and a lumped mass matrix. The same
identification applies term by term in \eqref{eq:pinoac}, but it now touches \emph{two} terms: the
time derivative and the reaction both lose their $M$-weighting, not just the source. So the PINO
Allen--Cahn residual is $h^{-2}$ times a $P_1$, mass-lumped version of \eqref{eq:fem} --- a
consistent scheme, but a different one from the $Q_1$ consistent-mass scheme the reference
trajectory solves.

This matters more than in the Poisson case because the reference data is generated by a FEM Newton
solve of \eqref{eq:fem}. The PINO residual is therefore \emph{not} zero at the reference solution:
it is zero at the solution of its own scheme, which differs at $O(h^2)$. That is the exact analogue
of the discretisation-error floor we removed for Poisson by adding
\texttt{--dataset\_solution fem}, and no such switch exists on the Allen--Cahn side, because the
reference is a time-stepped Newton solve rather than a single linear solve.

\section{What carries over unchanged}

\begin{itemize}
  \item \textbf{Finite differences, not FFT.} The reference uses FFT for periodic problems and
  central differences for Dirichlet ones (Darcy). Allen--Cahn here is homogeneous Dirichlet on the
  unit square, so central differences is the faithful choice, as for Poisson.
  \item \textbf{Interior-only evaluation.} Boundary rows and columns are sliced off, exactly as in
  \texttt{FDM\_Darcy}.
  \item \textbf{Hard boundary conditions in the model.} The Poisson arm now zeroes boundary nodes
  rather than using the $\sin(\pi x)\sin(\pi y)$ mollifier, so that it differs from our
  \texttt{galerkin} arm only in the residual. The Allen--Cahn arm must do the same. Note the
  mollifier was never defensible here anyway: it encodes a decay profile with no relation to a
  phase-field interface.
  \item \textbf{\texttt{detach\_coupling} / pushforward.} \texttt{PINOACLoss} takes the same
  \texttt{detach\_coupling} flag as \texttt{ACGalerkinLoss}, set from \texttt{--bptt\_mode}, so the
  two arms agree on what the gradient flows through and differ only in the residual. All
  \texttt{realistic\_ac} arms use \texttt{pushforward}.
\end{itemize}

\section{PI-DeepONet, in one paragraph}

The same residual \eqref{eq:pinoac} with $\Delta_h$ replaced by the autodiff Laplacian through the
trunk's coordinate input. One structural difference: the rollout must evaluate the model at
\emph{every} node to feed the next frame back, whereas \texttt{PINOACLoss} averages over the
interior only --- so \texttt{PIDeepONetACLoss} carries an explicit \texttt{interior\_mask}. Without
it the two baselines' losses would differ by the ratio of node counts alone ($6\,\%$ at $65^2$,
$3\,\%$ at $128^2$) and their learning rates would not be comparable.

\section{Summary}

\begin{center}
\begin{tabular}{lll}
\toprule
& Poisson & Allen--Cahn \\
\midrule
residual        & $-\Delta_h u - f$        & $\frac{u^{n+1}-u^n}{\dt} - a^2\Delta_h u^{n+1} - \rho$ \\
derivative      & \multicolumn{2}{c}{central finite differences --- shared} \\
evaluated on    & \multicolumn{2}{c}{interior nodes only --- shared} \\
Dirichlet BC    & \multicolumn{2}{c}{boundary nodes zeroed in the model --- shared} \\
reduction       & relative, $\|\cdot\|/\|f\|$ & \textbf{mse} (no right-hand side exists) \\
loss geometry   & cone, non-smooth at 0     & \textbf{quadratic}, $\cond \sim \cond(J)^2$ \\
linearised op.  & $A$, $\cond=\Theta(h^{-2})$ & $M/\dt + a^2A + 3\epsilon^2M\!\operatorname{diag}(u^2)$ \\
dropping $M$    & lumped load               & lumped load \emph{and} lumped time derivative \\
zero at ref.?   & yes (with \texttt{fem} labels) & \textbf{no} --- differs at $O(h^2)$ \\
\bottomrule
\end{tabular}
\end{center}

\end{document}
